% !TeX root = ../main.tex
%------------------------------------------------------------------------------
%  Appendix E -- Typical Material Constants and Parameter Ranges
%  Values compiled from the standard references (Hoek & Brown; Brady & Brown;
%  Hoek, Kaiser & Bawden; Mair & Taylor). Every value must be verified against
%  its source before use in print.
%------------------------------------------------------------------------------
\chapter{Typical Material Constants and Parameter Ranges}
\label{app:constants}

\noindent\textbf{Tabulated values are for checking, not for designing.}\enspace
Every value in this appendix is a typical range drawn from published
compilations: \citet{hoekbrown1980underground}, \citet{bradybrown2006rock}
and \citet{hoek1995support} for rock properties and support elements;
\citet{marinos2000gsi} and \citet{hoek2002hb} for the Geological Strength
Index, the Hoek--Brown constants and the disturbance factor;
\citet{barton1977shear} for discontinuity friction;
\citet{mair1996settlement} for volume loss and trough width; and
\citet{burland1974settlement} for the damage categories. They are useful for sanity-checking a measured value, for
preliminary study before site data exists, and for setting sensitivity bounds.
They are not a substitute for project-specific testing, and a design relying
on a published typical range rather than on test data should say so
explicitly.

\section{Intact rock and Hoek--Brown constants}
\label{sec:intact-rock-and-hoek-brown-constants}

\paragraph{Strength and stiffness by rock type.}
Table~\ref{tab:intact-props} gives indicative ranges. The scatter within a
rock type routinely exceeds the difference between types, so the table
identifies an implausible value rather than supplying a design one.

\begin{table}[htbp]
\caption{Indicative intact rock properties. Anisotropic rocks
(schist, slate, shale) vary strongly with the angle between loading and
fabric; the ranges given span that variation.}
\label{tab:intact-props}
\centering
\small
\begin{tabularx}{\textwidth}{l c c c X}
\toprule
\textbf{Rock type} & \boldmath{$\sci$} \textbf{(MPa)} & \boldmath{$\Ei$}
\textbf{(GPa)} & \boldmath{$\nu$} & \textbf{Unit weight (kN/m\textsuperscript{3})} \\
\midrule
Granite            & 100--250 & 40--70 & 0.15--0.25 & 26--27 \\
Basalt             & 100--300 & 40--80 & 0.15--0.25 & 27--30 \\
Sandstone          & 20--170  & 5--40  & 0.15--0.30 & 22--26 \\
Limestone          & 30--250  & 20--70 & 0.15--0.30 & 24--27 \\
Siltstone          & 10--120  & 5--30  & 0.15--0.30 & 23--26 \\
Shale / mudstone   & 5--100   & 2--20  & 0.20--0.35 & 22--26 \\
Gneiss             & 80--250  & 30--70 & 0.15--0.25 & 26--29 \\
Schist             & 20--150  & 5--50  & 0.15--0.30 & 25--28 \\
Slate              & 60--250  & 20--60 & 0.15--0.25 & 26--28 \\
\bottomrule
\end{tabularx}
\end{table}

The modulus ratio $MR = \Ei/\sci$ is often more stable than either quantity
alone, and where $\Ei$ has not been measured it may be estimated as
$\Ei = MR\,\sci$. The ranges of Table~\ref{tab:intact-props}, taken end for
end, imply ratios of roughly 200 to 400 for shale and mudstone, 230 to 250
for sandstone, 280 to 670 for limestone, and 280 to 400 for granite and
gneiss; published compilations by rock type give comparable
figures~\citep{hoekdiederichs2006}.

\paragraph{Hoek--Brown $\mi$ values.}
The constant $\mi$ describes the curvature of the intact failure envelope and
is grouped by rock type and texture, as in Table~\ref{tab:mi}. Values are
approximate and are conventionally quoted with an uncertainty of a few units.

\begin{table}[htbp]
\caption{Hoek--Brown constant $\mi$ for intact rock. Anisotropic rocks are
quoted for loading normal to the fabric; values fall substantially for loading
at an unfavourable angle.}
\label{tab:mi}
\centering
\small
\begin{tabularx}{\textwidth}{X c X c}
\toprule
\textbf{Rock type} & \boldmath{$\mi$} & \textbf{Rock type} & \boldmath{$\mi$} \\
\midrule
Granite            & 32 & Sandstone          & 17 \\
Granodiorite       & 29 & Siltstone          & 7 \\
Diorite            & 25 & Claystone          & 4 \\
Gabbro             & 27 & Greywacke          & 18 \\
Basalt             & 25 & Limestone, crystalline & 12 \\
Rhyolite           & 25 & Limestone, micritic    & 9 \\
Andesite           & 25 & Dolomite           & 9 \\
Gneiss             & 28 & Marble             & 9 \\
Schist             & 12 & Quartzite          & 20 \\
Phyllite           & 7  & Slate              & 7 \\
Amphibolite        & 26 & Chalk              & 7 \\
\bottomrule
\end{tabularx}
\end{table}

Interpolation for rocks not listed should be made by texture and mineralogy
rather than by name: coarse-grained crystalline rocks tend to high $\mi$, and
fine-grained clay-rich rocks to low.

\section{Rock mass parameters}
\label{sec:rock-mass-parameters-app}

\paragraph{Geological Strength Index.}
$\GSI$ is estimated from rock mass structure and discontinuity surface
condition, as in Table~\ref{tab:gsi}. It is deliberately a range rather than a
number, and a $\GSI$ quoted to a single unit misrepresents the evidence.

\begin{table}[htbp]
\caption{Indicative $\GSI$ ranges by structure and surface condition. The five surface classes of Chapter~\ref{ch:classification} are collapsed to three, good spanning very good and good and poor spanning poor and very poor; a dash marks a combination that does not occur.}
\label{tab:gsi}
\centering
\small
\begin{tabularx}{\textwidth}{X c c c}
\toprule
\textbf{Structure} & \textbf{Good} & \textbf{Fair} & \textbf{Poor} \\
\midrule
Intact or massive        & 80--90 & --     & --     \\
Blocky                   & 65--80 & 55--70 & 40--55 \\
Very blocky              & 55--70 & 45--60 & 30--45 \\
Blocky/disturbed, seamy  & 45--60 & 35--50 & 25--40 \\
Disintegrated            & 30--45 & 25--40 & 15--30 \\
Laminated, sheared       & --     & 15--30 & 10--20 \\
\bottomrule
\end{tabularx}
\end{table}

Where classification data exist, the correlations
$\GSI \approx \RMR'_{89} - 5$ (for $\RMR'_{89} > 23$) and
$\GSI \approx 9\ln Q' + 44$ provide a cross-check~\citep{hoek1995support}. Here $\RMR'_{89}$ is the
rating of Chapter~\ref{ch:classification} taken dry and without orientation
adjustment, and $Q'$ is the modified quality index
$(\RQD/J_n)(J_r/J_a)$, that is $Q$ with $J_w$ and SRF set to unity. They agree with chart
estimation to within about five points in competent rock and diverge in poor
ground, where the chart is preferred.

\paragraph{Disturbance factor.}
Table~\ref{tab:d-factor} gives $\Dfac$ by excavation method. It applies to the
damaged annulus only, not to the whole rock mass, and applying $\Dfac$
throughout a numerical model is a common and severely conservative error. The
thickness over which it applies depends on the excavation class. For blasted
tunnels a practical rule is to apply it over a thickness comparable to the
blast damage zone --- typically one to two metres --- and to use $\Dfac = 0$
beyond. For the slope rows the damaged zone scales with the bench or slope
height rather than being a fixed thickness. The value for mechanical
excavation in squeezing ground represents strain-induced disturbance, which
extends with the plastic zone rather than over a fixed annulus.

\begin{table}[htbp]
\caption{Disturbance factor by excavation method.}
\label{tab:d-factor}
\centering
\small
\begin{tabularx}{\textwidth}{X c}
\toprule
\textbf{Excavation method} & \boldmath{$\Dfac$} \\
\midrule
Mechanical excavation, minimal disturbance & 0 \\
Good-quality controlled blasting, smooth wall & 0.5 \\
Poor-quality production blasting & 0.8 \\
Very poor blasting, heavy damage & 1.0 \\
Mechanical excavation in squeezing ground & 0--0.5 \\
Slope excavation, small-scale controlled blasting & 0.7 \\
Slope excavation, production blasting & 1.0 \\
\bottomrule
\end{tabularx}
\end{table}

\section{Discontinuities and soils}
\label{sec:discontinuities-and-soils}

\paragraph{Discontinuity strength.}
Table~\ref{tab:joint} gives basic friction angles $\phi_b$ for clean,
unweathered surfaces. For a rough unfilled surface the mobilised friction
angle is
\[
\phi_b + \mathrm{JRC}\,\log_{10}\!\left(\frac{\mathrm{JCS}}{\sigma_n'}\right),
\]
in which JRC is the joint roughness coefficient, JCS the joint wall
compressive strength and $\sigma_n'$ the effective normal
stress~\citep{barton1977shear}. The roughness contribution is therefore not
a constant: it falls as the normal stress rises and vanishes as $\sigma_n'$
approaches JCS, which is the mechanism described in
Chapter~\ref{ch:engineering-geology}. Weathering and infill reduce the
angle further.

\begin{table}[htbp]
\caption{Basic friction angle $\phi_b$ for clean unweathered discontinuity
surfaces. Residual friction angles for weathered surfaces are typically
several degrees lower.}
\label{tab:joint}
\centering
\small
\begin{tabularx}{\textwidth}{X c X c}
\toprule
\textbf{Rock} & \boldmath{$\phi_b$} \textbf{(\textdegree)} & \textbf{Rock} & \boldmath{$\phi_b$} \textbf{(\textdegree)} \\
\midrule
Granite   & 29--35 & Sandstone & 26--35 \\
Basalt    & 31--38 & Siltstone & 27--31 \\
Limestone & 31--37 & Shale     & 27--32 \\
Dolomite  & 27--35 & Slate     & 25--30 \\
Gneiss    & 26--29 & Chalk     & 30--41 \\
\bottomrule
\end{tabularx}
\end{table}

The joint roughness coefficient ranges from 0--2 for smooth planar surfaces to
18--20 for very rough undulating ones, and it is scale-dependent: a value
measured on a 100\,mm profile overestimates the roughness of a joint several
metres across. Joint wall compressive strength equals $\sci$ for fresh
surfaces and falls to a quarter of it or less for weathered ones.

Where infill thickness exceeds the asperity amplitude, the surfaces are
decoupled and the strength is that of the infill: typically $\phi' = 10$ to
$20^\circ$ for clay gouge, with little or no effective cohesion.

\paragraph{Soils for tunnelling.}
Table~\ref{tab:soil-params} gives indicative ranges for soils commonly
encountered in urban tunnelling. The drained parameters apply to long-term
conditions. Face stability in clay during excavation is a total-stress check
on the undrained shear strength $s_u$, through the stability ratio of
Section~\ref{sec:face-stability-analysis}.

\begin{table}[htbp]
\caption{Indicative soil parameters: drained $\phi'$, $c'$ and $E'$, and
undrained shear strength $s_u$ for clays. Small-strain stiffness governing
ground movement outside the yield zone is commonly three to five times the
values of $E'$ tabulated. Soft clay is taken below about 40\,kPa and stiff
clay above about 75\,kPa, firm clay lying between; the soft and stiff clay
rows of Table~\ref{tab:volume-loss} are read with the same division.}
\label{tab:soil-params}
\centering
\small
\begin{tabularx}{\textwidth}{X c c c c}
\toprule
\textbf{Soil} & \boldmath{$\phi'$} \textbf{(\textdegree)} & \boldmath{$c'$}
\textbf{(kPa)} & \boldmath{$E'$} \textbf{(MPa)} & \boldmath{$s_u$} \textbf{(kPa)} \\
\midrule
Loose sand            & 28--32 & 0      & 10--25  & --- \\
Dense sand            & 35--45 & 0      & 50--150 & --- \\
Gravel                & 35--45 & 0      & 60--200 & --- \\
Soft clay             & 18--25 & 0--10  & 2--10   & 20--40 \\
Stiff clay            & 20--28 & 5--30  & 20--60  & 75--150 \\
Hard clay / residual  & 24--32 & 10--50 & 40--150 & above 150 \\
\bottomrule
\end{tabularx}
\end{table}

\section{Support elements and concrete}
\label{sec:support-elements-and-concrete}

\paragraph{Support element capacities.}
Table~\ref{tab:support-props} gives representative capacities.
Section~\ref{sec:grc-scc-appendix} gives stiffness and capacity expressions
for concrete and sprayed concrete rings, segmental rings, steel sets and
anchored or grouted bolts. It does not treat friction and energy-absorbing
bolts, cables or mesh, whose stiffness is taken from the load--displacement
response measured in pull-out tests or, for mesh, panel tests. For the
steel-set and bolt expressions, the modulus of steel is 200 to
210~GPa for sets and bars alike; the structural sections used as sets have
yield stresses of 275 to 355~MPa and cross-sectional areas of roughly 30 to
80~cm\textsuperscript{2}, and are placed at longitudinal spacings of 0.75 to
1.5~m.

\begin{table}[htbp]
\caption{Representative support element capacities.}
\label{tab:support-props}
\centering
\small
\begin{tabularx}{\textwidth}{X c X}
\toprule
\textbf{Element} & \textbf{Capacity} & \textbf{Remarks} \\
\midrule
Grouted rebar, 20\,mm       & 130--180\,kN & Bar yield load; see text \\
Grouted rebar, 25\,mm       & 200--280\,kN & Bar yield load; see text \\
Friction bolt (split set)   & 40--100\,kN  & Low capacity, high deformability \\
Expandable friction bolt    & 100--200\,kN & \\
Cable, single strand        & 200--260\,kN & \\
Cable, multi-strand anchor  & 500--2000\,kN & Pre-tensioned; long free length \\
Energy-absorbing bolt       & 100--200\,kN & Rated in kJ, not kN; 20--50\,kJ typical \\
Welded mesh                 & 20--40\,kN/m & Brittle failure at welds \\
Chain-link mesh             & 30--60\,kN/m & Large deformation capacity \\
\bottomrule
\end{tabularx}
\end{table}

The rebar values are the yield loads of the bar itself, and they scale with bar
area rather than perimeter. The capacity of an installed bar may be lower: it
depends on the embedment length, the grout quality and the borehole condition,
and where the embedment is short the bond governs before the bar yields. A
bond capacity is meaningful only per metre of embedment and is not tabulated
here.

\paragraph{Concrete and shotcrete.}
Sprayed concrete is normally specified at 30 to 45\,MPa at 28 days, with
accelerator dosage reducing long-term strength by 15 to 30\,\% relative to the
unaccelerated mix --- an allowance that belongs in the design strength rather
than in the test results.

Early strength governs the excavation cycle. Indicative development for a
wet-mix accelerated shotcrete is 0.2 to 0.5\,MPa at one hour, 1 to 3\,MPa at
three hours, 5 to 10\,MPa at twelve hours and 15 to 25\,MPa at one day. Elastic
modulus rises with strength, reaching roughly 10\,GPa at one day and 25 to
30\,GPa at 28 days, and the creep coefficient of young shotcrete is large
enough --- of order 2 to 3 --- to relax a substantial part of any load applied
in the first days.

Precast segments are typically specified at 45 to 60\,MPa, with steel fibre
dosages of 30 to 45\,kg/m\textsuperscript{3} and residual flexural strength
rather than dosage as the specified performance.

\section{Settlement parameters and conversions}
\label{sec:settlement-parameters-and-conversions}

\paragraph{Volume loss and trough width.}
Table~\ref{tab:volume-loss} gives volume loss achievable with good
construction control. Poor control, mixed face conditions or inadequate
grouting can multiply these by a factor of two or more.

\begin{table}[htbp]
\caption{Indicative volume loss and trough width parameter. The trough width
$i = K z_0$, where $z_0$ is the depth to tunnel axis; the trough and $V_L$ are
defined in Section~\ref{sec:gaussian-trough}.}
\label{tab:volume-loss}
\centering
\small
\begin{tabularx}{\textwidth}{X c c}
\toprule
\textbf{Method and ground} & \boldmath{$V_L$} \textbf{(\%)} & \boldmath{$K$} \\
\midrule
Closed-face machine, stiff clay      & 0.5--1.5 & 0.4--0.5 \\
Closed-face machine, soft clay       & 1.0--2.0 & 0.5 \\
Closed-face machine, sand below water& 0.5--1.5 & 0.25--0.35 \\
Sprayed concrete lining, stiff clay  & 0.5--1.5 & 0.4--0.5 \\
Sprayed concrete lining, weak rock   & 0.3--1.0 & 0.3--0.4 \\
Open face, stiff clay                & 1.0--3.0 & 0.5 \\
\bottomrule
\end{tabularx}
\end{table}

Damage category thresholds in terms of limiting tensile strain are
conventionally: negligible below 0.05\,\%; very slight 0.05--0.075\,\%; slight
0.075--0.15\,\%; moderate 0.15--0.3\,\%; and severe to very severe above 0.3\,\%, where the
two highest categories are distinguished by the visible damage rather than by
strain.

\paragraph{Conversions.}
$1\,\mathrm{MPa} = 10.2\,\mathrm{kgf/cm^2} = 145\,\mathrm{psi}$;
$1\,\mathrm{MN/m^2} = 1\,\mathrm{MPa}$;
$1\,\mathrm{bar} = 100\,\mathrm{kPa}$;
a head of $1\,\mathrm{m}$ of water $= 9.81\,\mathrm{kPa}$;
$1\,\mathrm{Lugeon} \approx 1.3\times10^{-7}\,\mathrm{m/s}$;
$1\,\mathrm{tonne\,force} = 9.81\,\mathrm{kN}$.

\chapterreferences
